% !TeX program = lualatex
% !TeX root = predavanje06.tex
% ============================================================
% Zajednički preamble za sva predavanja RWA kolegija
% Uključiti s: % ============================================================
% Zajednički preamble za sva predavanja RWA kolegija
% Uključiti s: % ============================================================
% Zajednički preamble za sva predavanja RWA kolegija
% Uključiti s: \input{../common/preamble.tex}
% ============================================================

\documentclass[aspectratio=149,11pt]{beamer}
\usetheme{CambridgeUS}

\usepackage{fontspec}              % radi samo sa XeLaTeX/LuaLaTeX
\usepackage[croatian]{babel}
\usepackage{csquotes}              % \enquote{} → hrvatski navodnici „..."
\usepackage{booktabs}
\usepackage{array}
\usepackage{listings}
\usepackage{xcolor}
\usepackage{colortbl}
\usepackage{hyperref}
\usepackage{tikz}
\usetikzlibrary{positioning,arrows.meta}

% ---------- Tipografija: IBM Plex ----------
\setsansfont{IBMPlexSans}[
  Extension      = .otf,
  UprightFont    = *-Regular,
  ItalicFont     = *-Italic,
  BoldFont       = *-SemiBold,
  BoldItalicFont = *-SemiBoldItalic
]
\setmainfont{IBMPlexSans}[
  Extension      = .otf,
  UprightFont    = *-Regular,
  ItalicFont     = *-Italic,
  BoldFont       = *-SemiBold,
  BoldItalicFont = *-SemiBoldItalic
]
\setmonofont{IBMPlexMono}[
  Extension      = .otf,
  UprightFont    = *-Regular,
  ItalicFont     = *-Italic,
  BoldFont       = *-SemiBold,
  BoldItalicFont = *-SemiBoldItalic,
  Scale          = 0.92
]

% ---------- Paleta ----------
\definecolor{ink}{HTML}{0C4685}      % primarna tamnoplava
\definecolor{ink2}{HTML}{6E7BAA}     % sekundarna plavo-siva
\definecolor{tint}{HTML}{EAEEF6}     % svjetli plavi tint
\definecolor{textclr}{HTML}{1A222E}  % glavni tekst
\definecolor{mutedclr}{HTML}{6B7484} % podnožje, sekundarni tekst
\definecolor{bgclr}{HTML}{F6F7F9}    % pozadina slajda
\definecolor{lineclr}{HTML}{D9DEE7}  % tanke linije tablica
\definecolor{divrule}{HTML}{9DB4D6}  % akcent na tamnim sekcijskim slajdovima
\definecolor{divtext}{HTML}{E6ECF6}  % tekst na tamnim sekcijskim slajdovima

\setbeamercolor{normal text}{fg=textclr,bg=bgclr}
\setbeamercolor{structure}{fg=ink}
\setbeamercolor{alerted text}{fg=ink}

\setbeamercolor{title}{fg=ink,bg=tint}
\setbeamercolor{subtitle}{fg=mutedclr}
\setbeamercolor{author}{fg=textclr}
\setbeamercolor{date}{fg=mutedclr}
\setbeamercolor{frametitle}{fg=ink,bg=tint}

\setbeamercolor{section in head/foot}{fg=white,bg=ink}
\setbeamercolor{subsection in head/foot}{fg=ink,bg=tint}
\setbeamercolor{author in head/foot}{fg=white,bg=ink}
\setbeamercolor{title in head/foot}{fg=textclr,bg=tint}
\setbeamercolor{date in head/foot}{fg=mutedclr,bg=tint}

\setbeamercolor{block title}{fg=white,bg=ink}
\setbeamercolor{block body}{fg=textclr,bg=tint}
\setbeamercolor{block title example}{fg=white,bg=ink2}
\setbeamercolor{block body example}{fg=textclr,bg=tint}
\setbeamercolor{itemize item}{fg=ink2}
\setbeamercolor{itemize subitem}{fg=ink2}
\setbeamercolor{enumerate item}{fg=ink2}

\setbeamertemplate{itemize items}[square]
\setbeamertemplate{navigation symbols}{}
\arrayrulecolor{lineclr}

% ---------- Kolegij, autori i logo ----------
\newcommand{\kolegij}{Programiranje za web}
\author{doc.\,dr.\,sc. Tomislav Petković \and mag.\,ing. Josip Torić}
\newcommand{\znak}[2]{\raisebox{-0.5\height}{\includegraphics[height=#1]{#2}}}
\titlegraphic{\znak{2.1cm}{../common/SIT_logo_tekst1.pdf}\hspace{1.8em}\znak{1.8cm}{../common/UNIZD_logo.pdf}}

% ---------- Zaglavlje ----------
% Jedna puna traka: lijevo sekcija, desno podsekcija.
\setbeamertemplate{headline}{%
  \leavevmode%
  \hbox{%
    \begin{beamercolorbox}[wd=\paperwidth,ht=2.25ex,dp=1ex,leftskip=.3cm,rightskip=.3cm]{section in head/foot}%
      \usebeamerfont{section in head/foot}\insertsectionhead\hfill\insertsubsectionhead%
    \end{beamercolorbox}%
  }%
}

% ---------- Podnožje ----------
% Lijevo: naziv trenutne prezentacije (kratki oblik naslova iz \title[...]{...})
% Sredina: naziv kolegija.  Desno: datum i redni broj slajda.
\setbeamertemplate{footline}{%
  \leavevmode%
  \hbox{%
    \begin{beamercolorbox}[wd=.34\paperwidth,ht=2.25ex,dp=1ex,center]{author in head/foot}%
      \usebeamerfont{author in head/foot}\insertshorttitle%
    \end{beamercolorbox}%
    \begin{beamercolorbox}[wd=.33\paperwidth,ht=2.25ex,dp=1ex,center]{title in head/foot}%
      \usebeamerfont{title in head/foot}\kolegij%
    \end{beamercolorbox}%
    \begin{beamercolorbox}[wd=.33\paperwidth,ht=2.25ex,dp=1ex,center]{date in head/foot}%
      \usebeamerfont{date in head/foot}\insertshortdate{}\hspace*{2em}%
      \insertframenumber{}\,/\,\inserttotalframenumber\hspace*{1ex}%
    \end{beamercolorbox}%
  }%
  \vskip0pt%
}

% ---------- Sekcijski slajd ----------
% \sekcija{Naslov sekcije}{jedna linija koja nabraja teme sekcije}
\newcommand{\sekcija}[2]{%
  \section{#1}%
  {%
    \setbeamercolor{background canvas}{bg=ink}%
    \setbeamertemplate{headline}{}%
    \begin{frame}
      \vfill
      \hspace*{0.03\textwidth}{\color{divrule}\rule{2.6cm}{1.1pt}}\par
      \vskip 0.7em
      \hspace*{0.03\textwidth}{\Large\bfseries\color{white}#1}\par
      \vskip 0.6em
      \hspace*{0.03\textwidth}\begin{minipage}{0.9\textwidth}\color{divtext}\small #2\end{minipage}\par
      \vfill
    \end{frame}%
  }%
}

% ---------- Listings ----------
\lstdefinelanguage{bash}{
  morekeywords={sudo,apt,brew,pip,pipx,uv,winget,python,git,flask,source,deactivate,code,gunicorn,docker,compose,heroku,render,railway,node,npm,pnpm,yarn,uvicorn,alembic,pytest,ruff,black,precommit},
  sensitive=false,
  morecomment=[l]{\#}
}
\lstdefinelanguage{python}{
  morekeywords={from,import,def,class,return,if,elif,else,for,in,while,try,except,raise,with,as,not,and,or,is,lambda,yield,True,False,None,await,async},
  sensitive=true,
  morecomment=[l]{\#},
  morestring=[b]',
  morestring=[b]"
}
\lstdefinelanguage{sql}{
  morekeywords={SELECT,FROM,WHERE,INSERT,INTO,VALUES,UPDATE,SET,DELETE,CREATE,TABLE,PRIMARY,KEY,FOREIGN,REFERENCES,UNIQUE,NOT,NULL,INDEX,ON,ALTER,DROP,CONSTRAINT},
  sensitive=false,
  morecomment=[l]{--}
}
\lstdefinelanguage{json}{
  morecomment=[l]{//},
  morestring=[b]",
  showstringspaces=false
}
\lstdefinelanguage{typescript}{
  morekeywords={import,export,from,const,let,var,function,return,if,else,for,of,in,while,switch,case,break,async,await,try,catch,finally,throw,new,class,this,extends,implements,interface,type,enum,readonly,public,private,protected,static,get,set,true,false,null,undefined,void,string,number,boolean,any,unknown,never,typeof,keyof,as},
  sensitive=true,
  morecomment=[l]{//},
  morecomment=[s]{/*}{*/},
  morestring=[b]',
  morestring=[b]",
  morestring=[b]`
}
\lstalias{javascript}{typescript}
\lstset{
  backgroundcolor=\color{tint},
  basicstyle=\ttfamily\footnotesize,
  keywordstyle=\color{ink}\bfseries,
  commentstyle=\color{mutedclr},
  frame=single,
  rulecolor=\color{lineclr},
  frameround=tttt,
  breaklines=true,
  tabsize=2,
  showstringspaces=false,
  upquote=true,
  columns=fullflexible,
  escapeinside={<@}{@>}
}
% ---------- Dijakritici u isječcima koda ----------
% listings čita bajt po bajt, pa višebajtni znak premjesti: "početni" ispadne
% "čpoetni". literate to ne popravlja. Rješenje je izlaz u LaTeX: sve između
% <@ i @> ne prolazi kroz listings, nego se slaže kao običan tekst, a unutar
% okvira je font ionako ttfamily, pa se ništa ne razlikuje.
%
%   async def get(...):
%       club = await repo.get_by_id(db, club_id)   <@\kom{dohvat kluba}@>
%       raise AppError("not_found", <@"Klub nije pronađen"@>, 404)
%
% Graničnici <@ i @> odabrani su jer se ne pojavljuju ni u jednom jeziku
% koji se na predavanjima citira; (* bi se sudaralo s def f(*args).
\newcommand{\kom}[1]{{\color{mutedclr}\# #1}}
% Isto za TypeScript, Vue i CSS, kojima komentar počinje s //:
%   const lifters = useLiftersStore()   <@\komts{dijeljeni popis}@>
\newcommand{\komts}[1]{{\color{mutedclr}// #1}}

% ---------- Činjenica naspram konvencije ----------
% Sve što nije propisano standardom ili jezikom, nego je dogovor struke,
% ide u ovaj okvir, da student zna koje je vrste tvrdnja pred njim.
% \begin{konvencija} ... \end{konvencija}
% \begin{konvencija}[Konvencija: nazivi putanja] ... \end{konvencija}
\newenvironment{konvencija}[1][Konvencija]
  {\begin{exampleblock}{#1}}
  {\end{exampleblock}}

% ---------- Helper macros ----------
% Stupac tablice zadane širine, poravnat ulijevo (bez rastezanja i dijeljenja riječi).
\newcolumntype{L}[1]{>{\raggedright\arraybackslash}p{#1}}
\newcommand{\codefile}[1]{\texttt{#1}}
\newcommand{\kbd}[1]{\texttt{#1}}

% ---------- Code frame environment ----------
% Automatski dodaje fragile i shrink=5, koristi se umjesto \begin{frame}
% Primjer: \begin{codeframe}{Naslov} ... \end{codeframe}
% Primjer s opcijama: \begin{codeframe}[shrink=15]{Naslov} ... \end{codeframe}
\newenvironment{codeframe}[2][]
  {\begin{frame}[fragile,shrink=5,#1]{#2}}
  {\end{frame}}

% ============================================================

\documentclass[aspectratio=149,11pt]{beamer}
\usetheme{CambridgeUS}

\usepackage{fontspec}              % radi samo sa XeLaTeX/LuaLaTeX
\usepackage[croatian]{babel}
\usepackage{csquotes}              % \enquote{} → hrvatski navodnici „..."
\usepackage{booktabs}
\usepackage{array}
\usepackage{listings}
\usepackage{xcolor}
\usepackage{colortbl}
\usepackage{hyperref}
\usepackage{tikz}
\usetikzlibrary{positioning,arrows.meta}

% ---------- Tipografija: IBM Plex ----------
\setsansfont{IBMPlexSans}[
  Extension      = .otf,
  UprightFont    = *-Regular,
  ItalicFont     = *-Italic,
  BoldFont       = *-SemiBold,
  BoldItalicFont = *-SemiBoldItalic
]
\setmainfont{IBMPlexSans}[
  Extension      = .otf,
  UprightFont    = *-Regular,
  ItalicFont     = *-Italic,
  BoldFont       = *-SemiBold,
  BoldItalicFont = *-SemiBoldItalic
]
\setmonofont{IBMPlexMono}[
  Extension      = .otf,
  UprightFont    = *-Regular,
  ItalicFont     = *-Italic,
  BoldFont       = *-SemiBold,
  BoldItalicFont = *-SemiBoldItalic,
  Scale          = 0.92
]

% ---------- Paleta ----------
\definecolor{ink}{HTML}{0C4685}      % primarna tamnoplava
\definecolor{ink2}{HTML}{6E7BAA}     % sekundarna plavo-siva
\definecolor{tint}{HTML}{EAEEF6}     % svjetli plavi tint
\definecolor{textclr}{HTML}{1A222E}  % glavni tekst
\definecolor{mutedclr}{HTML}{6B7484} % podnožje, sekundarni tekst
\definecolor{bgclr}{HTML}{F6F7F9}    % pozadina slajda
\definecolor{lineclr}{HTML}{D9DEE7}  % tanke linije tablica
\definecolor{divrule}{HTML}{9DB4D6}  % akcent na tamnim sekcijskim slajdovima
\definecolor{divtext}{HTML}{E6ECF6}  % tekst na tamnim sekcijskim slajdovima

\setbeamercolor{normal text}{fg=textclr,bg=bgclr}
\setbeamercolor{structure}{fg=ink}
\setbeamercolor{alerted text}{fg=ink}

\setbeamercolor{title}{fg=ink,bg=tint}
\setbeamercolor{subtitle}{fg=mutedclr}
\setbeamercolor{author}{fg=textclr}
\setbeamercolor{date}{fg=mutedclr}
\setbeamercolor{frametitle}{fg=ink,bg=tint}

\setbeamercolor{section in head/foot}{fg=white,bg=ink}
\setbeamercolor{subsection in head/foot}{fg=ink,bg=tint}
\setbeamercolor{author in head/foot}{fg=white,bg=ink}
\setbeamercolor{title in head/foot}{fg=textclr,bg=tint}
\setbeamercolor{date in head/foot}{fg=mutedclr,bg=tint}

\setbeamercolor{block title}{fg=white,bg=ink}
\setbeamercolor{block body}{fg=textclr,bg=tint}
\setbeamercolor{block title example}{fg=white,bg=ink2}
\setbeamercolor{block body example}{fg=textclr,bg=tint}
\setbeamercolor{itemize item}{fg=ink2}
\setbeamercolor{itemize subitem}{fg=ink2}
\setbeamercolor{enumerate item}{fg=ink2}

\setbeamertemplate{itemize items}[square]
\setbeamertemplate{navigation symbols}{}
\arrayrulecolor{lineclr}

% ---------- Kolegij, autori i logo ----------
\newcommand{\kolegij}{Programiranje za web}
\author{doc.\,dr.\,sc. Tomislav Petković \and mag.\,ing. Josip Torić}
\newcommand{\znak}[2]{\raisebox{-0.5\height}{\includegraphics[height=#1]{#2}}}
\titlegraphic{\znak{2.1cm}{../common/SIT_logo_tekst1.pdf}\hspace{1.8em}\znak{1.8cm}{../common/UNIZD_logo.pdf}}

% ---------- Zaglavlje ----------
% Jedna puna traka: lijevo sekcija, desno podsekcija.
\setbeamertemplate{headline}{%
  \leavevmode%
  \hbox{%
    \begin{beamercolorbox}[wd=\paperwidth,ht=2.25ex,dp=1ex,leftskip=.3cm,rightskip=.3cm]{section in head/foot}%
      \usebeamerfont{section in head/foot}\insertsectionhead\hfill\insertsubsectionhead%
    \end{beamercolorbox}%
  }%
}

% ---------- Podnožje ----------
% Lijevo: naziv trenutne prezentacije (kratki oblik naslova iz \title[...]{...})
% Sredina: naziv kolegija.  Desno: datum i redni broj slajda.
\setbeamertemplate{footline}{%
  \leavevmode%
  \hbox{%
    \begin{beamercolorbox}[wd=.34\paperwidth,ht=2.25ex,dp=1ex,center]{author in head/foot}%
      \usebeamerfont{author in head/foot}\insertshorttitle%
    \end{beamercolorbox}%
    \begin{beamercolorbox}[wd=.33\paperwidth,ht=2.25ex,dp=1ex,center]{title in head/foot}%
      \usebeamerfont{title in head/foot}\kolegij%
    \end{beamercolorbox}%
    \begin{beamercolorbox}[wd=.33\paperwidth,ht=2.25ex,dp=1ex,center]{date in head/foot}%
      \usebeamerfont{date in head/foot}\insertshortdate{}\hspace*{2em}%
      \insertframenumber{}\,/\,\inserttotalframenumber\hspace*{1ex}%
    \end{beamercolorbox}%
  }%
  \vskip0pt%
}

% ---------- Sekcijski slajd ----------
% \sekcija{Naslov sekcije}{jedna linija koja nabraja teme sekcije}
\newcommand{\sekcija}[2]{%
  \section{#1}%
  {%
    \setbeamercolor{background canvas}{bg=ink}%
    \setbeamertemplate{headline}{}%
    \begin{frame}
      \vfill
      \hspace*{0.03\textwidth}{\color{divrule}\rule{2.6cm}{1.1pt}}\par
      \vskip 0.7em
      \hspace*{0.03\textwidth}{\Large\bfseries\color{white}#1}\par
      \vskip 0.6em
      \hspace*{0.03\textwidth}\begin{minipage}{0.9\textwidth}\color{divtext}\small #2\end{minipage}\par
      \vfill
    \end{frame}%
  }%
}

% ---------- Listings ----------
\lstdefinelanguage{bash}{
  morekeywords={sudo,apt,brew,pip,pipx,uv,winget,python,git,flask,source,deactivate,code,gunicorn,docker,compose,heroku,render,railway,node,npm,pnpm,yarn,uvicorn,alembic,pytest,ruff,black,precommit},
  sensitive=false,
  morecomment=[l]{\#}
}
\lstdefinelanguage{python}{
  morekeywords={from,import,def,class,return,if,elif,else,for,in,while,try,except,raise,with,as,not,and,or,is,lambda,yield,True,False,None,await,async},
  sensitive=true,
  morecomment=[l]{\#},
  morestring=[b]',
  morestring=[b]"
}
\lstdefinelanguage{sql}{
  morekeywords={SELECT,FROM,WHERE,INSERT,INTO,VALUES,UPDATE,SET,DELETE,CREATE,TABLE,PRIMARY,KEY,FOREIGN,REFERENCES,UNIQUE,NOT,NULL,INDEX,ON,ALTER,DROP,CONSTRAINT},
  sensitive=false,
  morecomment=[l]{--}
}
\lstdefinelanguage{json}{
  morecomment=[l]{//},
  morestring=[b]",
  showstringspaces=false
}
\lstdefinelanguage{typescript}{
  morekeywords={import,export,from,const,let,var,function,return,if,else,for,of,in,while,switch,case,break,async,await,try,catch,finally,throw,new,class,this,extends,implements,interface,type,enum,readonly,public,private,protected,static,get,set,true,false,null,undefined,void,string,number,boolean,any,unknown,never,typeof,keyof,as},
  sensitive=true,
  morecomment=[l]{//},
  morecomment=[s]{/*}{*/},
  morestring=[b]',
  morestring=[b]",
  morestring=[b]`
}
\lstalias{javascript}{typescript}
\lstset{
  backgroundcolor=\color{tint},
  basicstyle=\ttfamily\footnotesize,
  keywordstyle=\color{ink}\bfseries,
  commentstyle=\color{mutedclr},
  frame=single,
  rulecolor=\color{lineclr},
  frameround=tttt,
  breaklines=true,
  tabsize=2,
  showstringspaces=false,
  upquote=true,
  columns=fullflexible,
  escapeinside={<@}{@>}
}
% ---------- Dijakritici u isječcima koda ----------
% listings čita bajt po bajt, pa višebajtni znak premjesti: "početni" ispadne
% "čpoetni". literate to ne popravlja. Rješenje je izlaz u LaTeX: sve između
% <@ i @> ne prolazi kroz listings, nego se slaže kao običan tekst, a unutar
% okvira je font ionako ttfamily, pa se ništa ne razlikuje.
%
%   async def get(...):
%       club = await repo.get_by_id(db, club_id)   <@\kom{dohvat kluba}@>
%       raise AppError("not_found", <@"Klub nije pronađen"@>, 404)
%
% Graničnici <@ i @> odabrani su jer se ne pojavljuju ni u jednom jeziku
% koji se na predavanjima citira; (* bi se sudaralo s def f(*args).
\newcommand{\kom}[1]{{\color{mutedclr}\# #1}}
% Isto za TypeScript, Vue i CSS, kojima komentar počinje s //:
%   const lifters = useLiftersStore()   <@\komts{dijeljeni popis}@>
\newcommand{\komts}[1]{{\color{mutedclr}// #1}}

% ---------- Činjenica naspram konvencije ----------
% Sve što nije propisano standardom ili jezikom, nego je dogovor struke,
% ide u ovaj okvir, da student zna koje je vrste tvrdnja pred njim.
% \begin{konvencija} ... \end{konvencija}
% \begin{konvencija}[Konvencija: nazivi putanja] ... \end{konvencija}
\newenvironment{konvencija}[1][Konvencija]
  {\begin{exampleblock}{#1}}
  {\end{exampleblock}}

% ---------- Helper macros ----------
% Stupac tablice zadane širine, poravnat ulijevo (bez rastezanja i dijeljenja riječi).
\newcolumntype{L}[1]{>{\raggedright\arraybackslash}p{#1}}
\newcommand{\codefile}[1]{\texttt{#1}}
\newcommand{\kbd}[1]{\texttt{#1}}

% ---------- Code frame environment ----------
% Automatski dodaje fragile i shrink=5, koristi se umjesto \begin{frame}
% Primjer: \begin{codeframe}{Naslov} ... \end{codeframe}
% Primjer s opcijama: \begin{codeframe}[shrink=15]{Naslov} ... \end{codeframe}
\newenvironment{codeframe}[2][]
  {\begin{frame}[fragile,shrink=5,#1]{#2}}
  {\end{frame}}

% ============================================================

\documentclass[aspectratio=149,11pt]{beamer}
\usetheme{CambridgeUS}

\usepackage{fontspec}              % radi samo sa XeLaTeX/LuaLaTeX
\usepackage[croatian]{babel}
\usepackage{csquotes}              % \enquote{} → hrvatski navodnici „..."
\usepackage{booktabs}
\usepackage{array}
\usepackage{listings}
\usepackage{xcolor}
\usepackage{colortbl}
\usepackage{hyperref}
\usepackage{tikz}
\usetikzlibrary{positioning,arrows.meta}

% ---------- Tipografija: IBM Plex ----------
\setsansfont{IBMPlexSans}[
  Extension      = .otf,
  UprightFont    = *-Regular,
  ItalicFont     = *-Italic,
  BoldFont       = *-SemiBold,
  BoldItalicFont = *-SemiBoldItalic
]
\setmainfont{IBMPlexSans}[
  Extension      = .otf,
  UprightFont    = *-Regular,
  ItalicFont     = *-Italic,
  BoldFont       = *-SemiBold,
  BoldItalicFont = *-SemiBoldItalic
]
\setmonofont{IBMPlexMono}[
  Extension      = .otf,
  UprightFont    = *-Regular,
  ItalicFont     = *-Italic,
  BoldFont       = *-SemiBold,
  BoldItalicFont = *-SemiBoldItalic,
  Scale          = 0.92
]

% ---------- Paleta ----------
\definecolor{ink}{HTML}{0C4685}      % primarna tamnoplava
\definecolor{ink2}{HTML}{6E7BAA}     % sekundarna plavo-siva
\definecolor{tint}{HTML}{EAEEF6}     % svjetli plavi tint
\definecolor{textclr}{HTML}{1A222E}  % glavni tekst
\definecolor{mutedclr}{HTML}{6B7484} % podnožje, sekundarni tekst
\definecolor{bgclr}{HTML}{F6F7F9}    % pozadina slajda
\definecolor{lineclr}{HTML}{D9DEE7}  % tanke linije tablica
\definecolor{divrule}{HTML}{9DB4D6}  % akcent na tamnim sekcijskim slajdovima
\definecolor{divtext}{HTML}{E6ECF6}  % tekst na tamnim sekcijskim slajdovima

\setbeamercolor{normal text}{fg=textclr,bg=bgclr}
\setbeamercolor{structure}{fg=ink}
\setbeamercolor{alerted text}{fg=ink}

\setbeamercolor{title}{fg=ink,bg=tint}
\setbeamercolor{subtitle}{fg=mutedclr}
\setbeamercolor{author}{fg=textclr}
\setbeamercolor{date}{fg=mutedclr}
\setbeamercolor{frametitle}{fg=ink,bg=tint}

\setbeamercolor{section in head/foot}{fg=white,bg=ink}
\setbeamercolor{subsection in head/foot}{fg=ink,bg=tint}
\setbeamercolor{author in head/foot}{fg=white,bg=ink}
\setbeamercolor{title in head/foot}{fg=textclr,bg=tint}
\setbeamercolor{date in head/foot}{fg=mutedclr,bg=tint}

\setbeamercolor{block title}{fg=white,bg=ink}
\setbeamercolor{block body}{fg=textclr,bg=tint}
\setbeamercolor{block title example}{fg=white,bg=ink2}
\setbeamercolor{block body example}{fg=textclr,bg=tint}
\setbeamercolor{itemize item}{fg=ink2}
\setbeamercolor{itemize subitem}{fg=ink2}
\setbeamercolor{enumerate item}{fg=ink2}

\setbeamertemplate{itemize items}[square]
\setbeamertemplate{navigation symbols}{}
\arrayrulecolor{lineclr}

% ---------- Kolegij, autori i logo ----------
\newcommand{\kolegij}{Programiranje za web}
\author{doc.\,dr.\,sc. Tomislav Petković \and mag.\,ing. Josip Torić}
\newcommand{\znak}[2]{\raisebox{-0.5\height}{\includegraphics[height=#1]{#2}}}
\titlegraphic{\znak{2.1cm}{../common/SIT_logo_tekst1.pdf}\hspace{1.8em}\znak{1.8cm}{../common/UNIZD_logo.pdf}}

% ---------- Zaglavlje ----------
% Jedna puna traka: lijevo sekcija, desno podsekcija.
\setbeamertemplate{headline}{%
  \leavevmode%
  \hbox{%
    \begin{beamercolorbox}[wd=\paperwidth,ht=2.25ex,dp=1ex,leftskip=.3cm,rightskip=.3cm]{section in head/foot}%
      \usebeamerfont{section in head/foot}\insertsectionhead\hfill\insertsubsectionhead%
    \end{beamercolorbox}%
  }%
}

% ---------- Podnožje ----------
% Lijevo: naziv trenutne prezentacije (kratki oblik naslova iz \title[...]{...})
% Sredina: naziv kolegija.  Desno: datum i redni broj slajda.
\setbeamertemplate{footline}{%
  \leavevmode%
  \hbox{%
    \begin{beamercolorbox}[wd=.34\paperwidth,ht=2.25ex,dp=1ex,center]{author in head/foot}%
      \usebeamerfont{author in head/foot}\insertshorttitle%
    \end{beamercolorbox}%
    \begin{beamercolorbox}[wd=.33\paperwidth,ht=2.25ex,dp=1ex,center]{title in head/foot}%
      \usebeamerfont{title in head/foot}\kolegij%
    \end{beamercolorbox}%
    \begin{beamercolorbox}[wd=.33\paperwidth,ht=2.25ex,dp=1ex,center]{date in head/foot}%
      \usebeamerfont{date in head/foot}\insertshortdate{}\hspace*{2em}%
      \insertframenumber{}\,/\,\inserttotalframenumber\hspace*{1ex}%
    \end{beamercolorbox}%
  }%
  \vskip0pt%
}

% ---------- Sekcijski slajd ----------
% \sekcija{Naslov sekcije}{jedna linija koja nabraja teme sekcije}
\newcommand{\sekcija}[2]{%
  \section{#1}%
  {%
    \setbeamercolor{background canvas}{bg=ink}%
    \setbeamertemplate{headline}{}%
    \begin{frame}
      \vfill
      \hspace*{0.03\textwidth}{\color{divrule}\rule{2.6cm}{1.1pt}}\par
      \vskip 0.7em
      \hspace*{0.03\textwidth}{\Large\bfseries\color{white}#1}\par
      \vskip 0.6em
      \hspace*{0.03\textwidth}\begin{minipage}{0.9\textwidth}\color{divtext}\small #2\end{minipage}\par
      \vfill
    \end{frame}%
  }%
}

% ---------- Listings ----------
\lstdefinelanguage{bash}{
  morekeywords={sudo,apt,brew,pip,pipx,uv,winget,python,git,flask,source,deactivate,code,gunicorn,docker,compose,heroku,render,railway,node,npm,pnpm,yarn,uvicorn,alembic,pytest,ruff,black,precommit},
  sensitive=false,
  morecomment=[l]{\#}
}
\lstdefinelanguage{python}{
  morekeywords={from,import,def,class,return,if,elif,else,for,in,while,try,except,raise,with,as,not,and,or,is,lambda,yield,True,False,None,await,async},
  sensitive=true,
  morecomment=[l]{\#},
  morestring=[b]',
  morestring=[b]"
}
\lstdefinelanguage{sql}{
  morekeywords={SELECT,FROM,WHERE,INSERT,INTO,VALUES,UPDATE,SET,DELETE,CREATE,TABLE,PRIMARY,KEY,FOREIGN,REFERENCES,UNIQUE,NOT,NULL,INDEX,ON,ALTER,DROP,CONSTRAINT},
  sensitive=false,
  morecomment=[l]{--}
}
\lstdefinelanguage{json}{
  morecomment=[l]{//},
  morestring=[b]",
  showstringspaces=false
}
\lstdefinelanguage{typescript}{
  morekeywords={import,export,from,const,let,var,function,return,if,else,for,of,in,while,switch,case,break,async,await,try,catch,finally,throw,new,class,this,extends,implements,interface,type,enum,readonly,public,private,protected,static,get,set,true,false,null,undefined,void,string,number,boolean,any,unknown,never,typeof,keyof,as},
  sensitive=true,
  morecomment=[l]{//},
  morecomment=[s]{/*}{*/},
  morestring=[b]',
  morestring=[b]",
  morestring=[b]`
}
\lstalias{javascript}{typescript}
\lstset{
  backgroundcolor=\color{tint},
  basicstyle=\ttfamily\footnotesize,
  keywordstyle=\color{ink}\bfseries,
  commentstyle=\color{mutedclr},
  frame=single,
  rulecolor=\color{lineclr},
  frameround=tttt,
  breaklines=true,
  tabsize=2,
  showstringspaces=false,
  upquote=true,
  columns=fullflexible,
  escapeinside={<@}{@>}
}
% ---------- Dijakritici u isječcima koda ----------
% listings čita bajt po bajt, pa višebajtni znak premjesti: "početni" ispadne
% "čpoetni". literate to ne popravlja. Rješenje je izlaz u LaTeX: sve između
% <@ i @> ne prolazi kroz listings, nego se slaže kao običan tekst, a unutar
% okvira je font ionako ttfamily, pa se ništa ne razlikuje.
%
%   async def get(...):
%       club = await repo.get_by_id(db, club_id)   <@\kom{dohvat kluba}@>
%       raise AppError("not_found", <@"Klub nije pronađen"@>, 404)
%
% Graničnici <@ i @> odabrani su jer se ne pojavljuju ni u jednom jeziku
% koji se na predavanjima citira; (* bi se sudaralo s def f(*args).
\newcommand{\kom}[1]{{\color{mutedclr}\# #1}}
% Isto za TypeScript, Vue i CSS, kojima komentar počinje s //:
%   const lifters = useLiftersStore()   <@\komts{dijeljeni popis}@>
\newcommand{\komts}[1]{{\color{mutedclr}// #1}}

% ---------- Činjenica naspram konvencije ----------
% Sve što nije propisano standardom ili jezikom, nego je dogovor struke,
% ide u ovaj okvir, da student zna koje je vrste tvrdnja pred njim.
% \begin{konvencija} ... \end{konvencija}
% \begin{konvencija}[Konvencija: nazivi putanja] ... \end{konvencija}
\newenvironment{konvencija}[1][Konvencija]
  {\begin{exampleblock}{#1}}
  {\end{exampleblock}}

% ---------- Helper macros ----------
% Stupac tablice zadane širine, poravnat ulijevo (bez rastezanja i dijeljenja riječi).
\newcolumntype{L}[1]{>{\raggedright\arraybackslash}p{#1}}
\newcommand{\codefile}[1]{\texttt{#1}}
\newcommand{\kbd}[1]{\texttt{#1}}

% ---------- Code frame environment ----------
% Automatski dodaje fragile i shrink=5, koristi se umjesto \begin{frame}
% Primjer: \begin{codeframe}{Naslov} ... \end{codeframe}
% Primjer s opcijama: \begin{codeframe}[shrink=15]{Naslov} ... \end{codeframe}
\newenvironment{codeframe}[2][]
  {\begin{frame}[fragile,shrink=5,#1]{#2}}
  {\end{frame}}


\title[Predavanje 6: Poslovna pravila]{Programiranje za web \\ Predavanje 6: Poslovna pravila i rokovi}
\date{12.11.2026.}

\begin{document}

\frame{\titlepage}

% ------------------------------------------------------------
\begin{frame}{Plan predavanja}
\begin{enumerate}
  \item Pravila koja shema ne može provjeriti
  \item Stanje izvedeno iz vremena
  \item Prijelaz stanja umjesto brisanja
  \item Pogreška kao odgovor
  \item Testiranje pravila koja ovise o vremenu
\end{enumerate}

\vspace{0.6em}
Dosad smo pitali tko šalje zahtjev i smije li dirati taj zapis. Ovo predavanje dodaje treće pitanje: smije li se radnja izvesti u ovom trenutku.
\end{frame}

% ------------------------------------------------------------
\begin{frame}{Ishodi učenja}
Nakon ovog predavanja student može:

\begin{itemize}
  \item razlikovati provjeru ulaznih podataka od poslovnog pravila i reći gdje koja stoji
  \item izvesti stanje zapisa iz vremena umjesto da ga pohrani
  \item zapisati rok tako da njegovo značenje ne ovisi o poslužitelju koji ga čita
  \item napisati pravilo koje dopušta ili odbija radnju ovisno o fazi
  \item obrazložiti izbor između brisanja zapisa i promjene oznake stanja
  \item prevesti domensku pogrešku u odgovor s dosljednim kodom i porukom
  \item napisati test pravila koje ovisi o vremenu, bez čekanja
\end{itemize}
\end{frame}

% ------------------------------------------------------------
\begin{frame}{Domena primjera}
Natjecanje ima dva roka, preliminarni i konačni. Oni dijele vrijeme na tri razdoblja, a u svakom klub smije nešto drugo:

\vspace{0.5em}
\begin{center}
{\small
\setlength{\tabcolsep}{0pt}
\renewcommand{\arraystretch}{1.4}
\begin{tabular}{@{}L{4.2cm}L{7.0cm}@{}}
  \toprule
  \textbf{Razdoblje} & \textbf{Što klub smije} \\
  \midrule
  prije preliminarnog roka & prijaviti natjecatelja, mijenjati i odjaviti prijavu \\
  između dvaju rokova & mijenjati i odjaviti postojeću prijavu, ali ne prijaviti novoga \\
  nakon konačnog roka & ništa od navedenoga \\
  \bottomrule
\end{tabular}
}
\end{center}

\vspace{0.6em}
Popis natjecatelja koji organizator šalje sucima nastaje nakon konačnog roka. Sve što se tada promijeni, promijenilo bi i taj popis, pa izmjene prestaju.
\end{frame}

% ------------------------------------------------------------
\begin{frame}[fragile]{Zahtjev koji prolazi svaku dosadašnju provjeru (I)}
Natjecanje 7 održano je prošle subote. Klub 12 za njega sada šalje prijavu:

\begin{lstlisting}
POST /competitions/7/registrations
Authorization: Bearer <token kluba 12>

{"lifter_id": 31, "category": "83", "total": 500}
\end{lstlisting}

\vspace{0.5em}
Koji od mehanizama koje smo dosad obradili sprječava da ovaj zahtjev prođe?

\pause
\vspace{0.5em}
Svaki mehanizam koji smo dosad napisali propušta ovaj zahtjev:

\begin{itemize}
  \itemsep0.4em
  \item \codefile{routers}: \emph{token} je valjan i nije istekao, a uloga \texttt{club} smije stvarati prijave
  \item \codefile{schemas}: tijelo ima sva obvezna polja, a kategorija je na popisu
  \item \codefile{services}: natjecatelj 31 pripada klubu 12, pa provjera vlasništva prolazi
\end{itemize}
\end{frame}

% ------------------------------------------------------------
\begin{frame}{Zahtjev koji prolazi svaku dosadašnju provjeru (II)}
Konačni rok istekao je tri tjedna prije natjecanja, a popis koji je organizator poslao sucima nastao je iz prijava zatečenih u tom trenutku. Rezultati su odavno upisani.

\vspace{0.9em}
Prijava bi se svejedno upisala. Uloga opisuje pošiljatelja, vlasništvo opisuje odnos pošiljatelja i zapisa, a shema opisuje oblik tijela. Nijedno se od toga ne mijenja protekom vremena, pa nijedna od triju provjera ovaj zahtjev ne vidi kao problem.

\vspace{0.9em}
Provjeru koja ga odbija zovemo \textbf{poslovnim pravilom}. Ona gleda stanje aplikacije u trenutku zahtjeva, a ne sam zahtjev.
\end{frame}

% ============================================================
\sekcija{Pravila koja shema ne može provjeriti}{Razlika prema provjeri ulaznih podataka, mjesto pravila, primjer koji traži zapis iz baze, granica između sheme i pravila.}

% ------------------------------------------------------------
\begin{frame}{Provjera ulaza i poslovno pravilo}
Zahtjev za prijavom može biti odbijen iz dvaju različitih razloga. Razlikuju se po tome što je trebalo pogledati da bi se odbio:

\vspace{0.5em}
\begin{center}
{\small
\setlength{\tabcolsep}{4pt}
\renewcommand{\arraystretch}{1.4}
\begin{tabular}{@{}L{3.1cm}L{3.8cm}L{3.8cm}@{}}
  \toprule
  & \textbf{Provjera ulaza} & \textbf{Poslovno pravilo} \\
  \midrule
  Gleda           & samo poslani sadržaj & stanje aplikacije \\
  Mjesto          & \codefile{schemas}   & \codefile{services} \\
  Primjer         & kategorija izvan popisa & rok je prošao \\
  Isti zahtjev sutra & jednako odbijen   & može proći \\
  \bottomrule
\end{tabular}
}
\end{center}

\vspace{0.6em}
Najlakše ih razdvaja zadnji redak. Neispravna kategorija neispravna je i sutra i za godinu dana. Prijava nakon roka uredna je po obliku, a odbijena je zbog trenutka u kojem je stigla.
\end{frame}

% ------------------------------------------------------------
\begin{frame}{Mjesto poslovnog pravila}
Pravilo o roku ne može stajati u shemi, jer shema nema pristup bazi i ne zna koje je natjecanje u pitanju.

\vspace{0.5em}
Ne može stajati ni u ruti. Isti rok odlučuje o prijavi, o izmjeni i o odjavi, a to su tri rute, pa bi se ista provjera pisala tri puta i tri bi se puta mogla zaboraviti.

\vspace{0.5em}
Ostaje sloj \codefile{services}: on vidi i pošiljatelja i zapis dohvaćen iz baze i stoji ispred svakog upisa. Zato sve tri rute pozivaju funkciju koja pravilo provodi na jednom mjestu.

\vspace{0.6em}
\begin{konvencija}
Nijedan standard ni \emph{framework} ne propisuje gdje poslovna pravila stoje. Odluka da ih nosi sloj između rute i baze proizlazi iz podjele odgovornosti koju smo postavili na početku kolegija, a ne iz zahtjeva alata.
\end{konvencija}
\end{frame}

% ------------------------------------------------------------
\begin{frame}[fragile]{Primjer: kategorija i spol (I)}
Težinske kategorije IPF-a razlikuju se po spolu, a shema zna oba popisa:

\begin{lstlisting}[language=python]
# app/schemas/registration.py
MALE_CATEGORIES = ["59", "66", "74", "83", "93", "105", "120", "120+"]
FEMALE_CATEGORIES = ["47", "52", "57", "63", "69", "76", "84", "84+"]
CATEGORIES_BY_GENDER = {"M": MALE_CATEGORIES, "F": FEMALE_CATEGORIES}
class RegistrationCreate(BaseModel):
    lifter_id: int
    category: str

    @field_validator("category")
    @classmethod
    def valid_category(cls, v: str) -> str:
        return _known_category(v)
\end{lstlisting}

\vspace{0.3em}
\texttt{\_known\_category} odbija sve što nije ni na jednom od dvaju popisa, pa kategorija \texttt{"999"} ne prolazi dalje od sheme. Odgovor je \texttt{422}.
\end{frame}

% ------------------------------------------------------------
\begin{frame}[fragile]{Primjer: kategorija i spol (II)}
Shema ne zna kojem natjecatelju prijava pripada. Zahtjev nosi samo njegovu oznaku \texttt{lifter\_id}, a spol stoji u zapisu u bazi:

\begin{lstlisting}[language=python]
# app/services/registration_service.py
def _check_category(lifter: Lifter, category: str) -> None:
    if category not in CATEGORIES_BY_GENDER[lifter.gender]:
        raise AppError(
            "invalid_category",
            f"Kategorija {category} ne odgovara spolu natjecatelja",
            400,
        )
\end{lstlisting}

\vspace{0.3em}
Kategorija \texttt{"47"} za muškog natjecatelja prolazi shemu i pada tek na pravilu, nakon dohvata zapisa: \texttt{invalid\_category}, \texttt{400}.

\vspace{0.1em}
\texttt{422} znači da poslani sadržaj nije u redu, a \texttt{400} da je sadržaj u redu, ali radnja u ovom stanju nije moguća.
\end{frame}

% ------------------------------------------------------------
\begin{frame}[fragile]{Cjelovito i djelomično tijelo (I)}
Granicu najjasnije pokazuje pravilo o poretku dvaju rokova. Tijelo kojim se zapis stvara nosi oba, pa ih shema uspoređuje sama:

\begin{lstlisting}[language=python]
# app/schemas/competition.py
class CompetitionCreate(BaseModel):
    prelim_deadline: AwareDatetime
    final_deadline: AwareDatetime
    @field_validator("final_deadline")
    @classmethod
    def final_after_prelim(cls, v, info):
        prelim = info.data.get("prelim_deadline")
        if prelim and v <= prelim:
            raise ValueError(
                "final_deadline mora biti nakon prelim_deadline"
            )
        return v
\end{lstlisting}

\vspace{0.3em}
Oba su roka u istom tijelu, pa shema obrnuti poredak sama odbija, s \texttt{422}.
\end{frame}

% ------------------------------------------------------------
\begin{frame}[fragile]{Cjelovito i djelomično tijelo (II)}
Izmjena smije poslati samo \texttt{prelim\_deadline}. Usporedba tada treba \texttt{final\_deadline} iz zapisa, pa \texttt{CompetitionUpdate} taj validator nema. Pravilo se seli u sloj koji vidi i poslana polja i zatečeni zapis:

\begin{lstlisting}[language=python]
# app/services/competition_service.py
for field, value in body.model_dump(exclude_unset=True).items():
    setattr(competition, field, value)
if competition.prelim_deadline >= competition.final_deadline:
    raise AppError("invalid_deadlines",
                   "final_deadline mora biti nakon prelim_deadline", 400)
\end{lstlisting}

\vspace{0.3em}
Isto pravilo, dva mjesta i dva koda: \texttt{422} pri stvaranju, \texttt{400} pri izmjeni. Razlika nije u pravilu, nego u tome nosi li zahtjev cijeli zapis ili samo promjenu.
\end{frame}

% ------------------------------------------------------------
\begin{frame}[fragile]{Vježba: dva odbijanja iste kategorije}
Kao klub \texttt{behemot@plreg.hr}, uz \emph{Zimski kup} u fazi \texttt{OPEN}, šaljete prijavu Ivana Perića i mijenjate samo kategoriju. Njegov \texttt{club\_id} vraća \texttt{GET /auth/me}, a oznaku natjecatelja \texttt{GET /clubs/\{club\_id\}/lifters}:

\begin{lstlisting}
POST /competitions/{comp_id}/registrations
{"lifter_id": <oznaka natjecatelja>, "category": "...", "total": 500}
\end{lstlisting}

\vspace{0.25em}
\begin{enumerate}
  \itemsep0.15em
  \item Pošaljite kategoriju \texttt{"999"} i zabilježite status i \texttt{code}.
  \item Pošaljite kategoriju \texttt{"47"}, koja postoji, ali je ženska.
  \item U \codefile{app/schemas/registration.py} nađite redak koji je odbio prvi zahtjev.
  \item U \codefile{app/services/registration\_service.py} nađite redak koji je odbio drugi.
  \item Pošaljite kategoriju \texttt{"83"} i provjerite da prijava prolazi.
\end{enumerate}

\vspace{0.3em}
Zašto prva provjera ne treba bazu, a druga treba?
\end{frame}

% ============================================================
\sekcija{Stanje izvedeno iz vremena}{Pohranjeno ili izračunato razdoblje, rok kao trenutak na vremenskoj osi, matrica dopuštenih radnji i njezin prijepis u kod.}

% ------------------------------------------------------------
\begin{frame}{Stupac ili izračun}
Rokovi mijenjaju koje su radnje dopuštene, pa aplikacija pri svakom zahtjevu treba odgovor na jedno pitanje: u kojem je razdoblju zapis u ovom trenutku.

\vspace{0.7em}
Taj se podatak može dobiti na dva načina:

\vspace{0.3em}
\begin{itemize}
  \itemsep0.5em
  \item \textbf{Pohraniti.} Natjecanje dobiva još jedan stupac, a vrijednost mu mijenja program koji se pokreće po rasporedu, mimo obrade zahtjeva.
  \item \textbf{Izračunati.} Stupca nema. Pri svakom se zahtjevu trenutno vrijeme usporedi s dvama rokovima i iz toga slijedi razdoblje.
\end{itemize}

\vspace{0.7em}
Prvi način čita jedno polje, drugi izvodi dvije usporedbe. Razlikuju se po tome može li podatak biti zastario.
\end{frame}

% ------------------------------------------------------------
\begin{frame}[fragile]{Razdoblje kao pohranjeni podatak}
Stupac bi stajao uz ostala polja natjecanja:

\begin{lstlisting}[language=python]
# app/models/competition.py, kad bi se faza pohranjivala
phase: Mapped[str] = mapped_column(String(20), default="open")
\end{lstlisting}

\vspace{0.25em}
Vrijednost ne mijenja nijedan zahtjev, nego program pokrenut po rasporedu:

\begin{lstlisting}[language=python]
<@\kom{zadatak koji se pokreće svakih pet minuta}@>
async def refresh_phases(db):
    now = datetime.now(timezone.utc)
    for comp in await competition_repo.get_all(db, 1000, 0):
        if now >= comp.final_deadline:
            comp.phase = "closed"
        elif now >= comp.prelim_deadline:
            comp.phase = "prelim_passed"
    await db.commit()
\end{lstlisting}

\vspace{0.25em}
Do sljedećeg pokretanja stupac laže, pa prijave prolaze i nakon roka.
\end{frame}

% ------------------------------------------------------------
\begin{frame}[fragile]{Imenovana razdoblja}
\begin{lstlisting}[language=python]
# app/core/phases.py
class Phase(str, Enum):
    OPEN = "open"
    PRELIM_PASSED = "prelim_passed"
    CLOSED = "closed"
\end{lstlisting}

\vspace{0.35em}
\begin{center}
{\small
\setlength{\tabcolsep}{0pt}
\renewcommand{\arraystretch}{1.2}
\begin{tabular}{@{}L{3.6cm}L{7.4cm}@{}}
  \toprule
  \textbf{Faza} & \textbf{Vrijedi} \\
  \midrule
  \texttt{OPEN}          & prije preliminarnog roka \\
  \texttt{PRELIM\_PASSED} & između preliminarnog i konačnog roka \\
  \texttt{CLOSED}        & nakon konačnog roka \\
  \bottomrule
\end{tabular}
}
\end{center}

\vspace{0.5em}
Razdoblja su nabrojiv skup, pa usporedba glasi \texttt{phase == Phase.OPEN}, bez datuma razasutih po kodu.
\end{frame}

% ------------------------------------------------------------
\begin{frame}[fragile]{Funkcija koja fazu računa}
\begin{lstlisting}[language=python]
# app/core/phases.py
def get_competition_phase(competition: Competition) -> Phase:
    now = datetime.now(timezone.utc)
    prelim = competition.prelim_deadline
    final = competition.final_deadline

    if now < prelim:
        return Phase.OPEN
    if now < final:
        return Phase.PRELIM_PASSED
    return Phase.CLOSED
\end{lstlisting}

\vspace{0.3em}
Funkcija prima natjecanje i vraća fazu. Ne dira bazu, ne podiže iznimke i ne odlučuje što je dopušteno.
\end{frame}

% ------------------------------------------------------------
\begin{frame}{Izvedena vrijednost}
\textbf{Rokovi su podatak, a faza zaključak iz njih.} Zaključak se zato računa pri svakom zahtjevu, a u bazi ga nema.

\vspace{0.9em}
Isto vrijedi za svaku vrijednost koja slijedi iz drugih: dob iz datuma rođenja, ukupan iznos iz stavaka računa, preostalo mjesto iz kapaciteta i broja zauzetih.

\vspace{0.9em}
Pohranjena izvedena vrijednost mora se osvježiti pri svakoj promjeni podataka iz kojih se izvodi, pa se pohranjuje tek kad je izračun preskup da se ponavlja.
\end{frame}

% ------------------------------------------------------------
\begin{frame}{Rok kao trenutak na vremenskoj osi}
Zapis \enquote{14.02.2027. u 23:59} sam po sebi ne određuje trenutak. U Zadru je to jedan trenutak, u Londonu drugi, sat kasnije. Usporedba s takvim rokom dala bi različit ishod ovisno o tome gdje poslužitelj radi.

\vspace{0.7em}
Rok zato nosi oznaku vremenske zone, i to na trima mjestima:

\begin{itemize}
  \itemsep0.4em
  \item shema traži tip \texttt{AwareDatetime}, pa zahtjev s vrijednošću bez oznake zone dobiva \texttt{422}
  \item stupac je \texttt{DateTime(timezone=True)}, u PostgreSQL-u \texttt{timestamptz}
  \item usporedba se izvodi u zoni UTC, s \texttt{datetime.now(timezone.utc)}
\end{itemize}

\end{frame}

% ------------------------------------------------------------
\begin{frame}[fragile]{Oznaka zone u shemi}
\begin{lstlisting}[language=python]
# app/schemas/competition.py
class CompetitionCreate(BaseModel):
    prelim_deadline: AwareDatetime
    final_deadline: AwareDatetime
\end{lstlisting}

\vspace{0.3em}
\texttt{AwareDatetime} je \texttt{datetime} uz zahtjev da vrijednost nosi pomak od UTC-a. Vrijednost bez njega ne ulazi u aplikaciju:

\begin{lstlisting}
422  {"prelim_deadline": "2027-02-14T23:59:00"}
201  {"prelim_deadline": "2027-02-14T23:59:00+01:00"}
\end{lstlisting}

\vspace{0.2em}
Bez te provjere vrijednost bi stigla do baze, a baza je ne bi odbila, nego bi je protumačila u svojoj zoni. Pogreška bi se vidjela tek kao rok koji istječe u krivi sat.
\end{frame}

% ------------------------------------------------------------
\begin{frame}[fragile]{Oznaka zone u stupcu}
\begin{lstlisting}[language=python]
# app/models/competition.py
prelim_deadline: Mapped[dt.datetime] = mapped_column(
    DateTime(timezone=True), nullable=False
)
\end{lstlisting}

\vspace{0.25em}
U PostgreSQL-u je to \texttt{timestamptz}, a bez argumenta goli \texttt{timestamp}:

\vspace{0.25em}
\begin{center}
{\small
\setlength{\tabcolsep}{6pt}
\renewcommand{\arraystretch}{1.3}
\begin{tabular}{@{}L{4.0cm}L{3.9cm}L{2.6cm}@{}}
  \toprule
   & \texttt{timestamptz} & \texttt{timestamp} \\
  \midrule
  Upisano             & \texttt{23:59+01:00} & \texttt{23:59+01:00} \\
  Pohranjeno          & \texttt{22:59} UTC   & \texttt{23:59} \\
  Pročitano u Zadru   & \texttt{23:59+01:00} & \texttt{23:59} \\
  Pročitano u Londonu & \texttt{22:59+00:00} & \texttt{23:59} \\
  \bottomrule
\end{tabular}
}
\end{center}

\vspace{0.25em}
Prvi pamti trenutak, drugi znamenke sata: pomak pri upisu tiho otpada.
\end{frame}

% ------------------------------------------------------------
\begin{frame}[fragile]{Oznaka zone u usporedbi}
\begin{lstlisting}[language=python]
# app/core/phases.py
now = datetime.now(timezone.utc)
if now < prelim:
    return Phase.OPEN
\end{lstlisting}

\vspace{0.3em}
\texttt{datetime.now()} bez argumenta vraća lokalno vrijeme poslužitelja, bez oznake zone. Takva se vrijednost s rokom iz baze ne može usporediti:

\begin{lstlisting}[language=python]
TypeError: can't compare offset-naive and offset-aware datetimes
\end{lstlisting}

\vspace{0.2em}
Iznimku dobivamo samo kad je jedna strana s oznakom, a druga bez nje. Ako su obje bez oznake, usporedba prolazi i daje ishod po satu poslužitelja, pa se pogreška vidi tek kad se poslužitelj preseli u drugu zonu.
\end{frame}

% ------------------------------------------------------------
\begin{frame}{Trenutak isteka roka}
Uvjet je napisan kao \texttt{now < prelim}. U trenutku koji rok određuje taj uvjet više ne vrijedi, pa natjecanje tada već nije u fazi \texttt{OPEN}. Rok u 23:59 znači da zadnja prijava mora stići prije 23:59, a ne u 23:59.

\vspace{0.9em}
\begin{konvencija}[Izbor, ne pravilo]
Granica bi se mogla napisati i kao \texttt{now <= prelim}, pa bi i trenutak roka još pripadao otvorenoj fazi. Nijedan standard o tome ne odlučuje. Važno je da odabrano značenje jednako glasi u propisu i u kodu.
\end{konvencija}

\vspace{0.9em}
Granice su mjesta na kojima se najlakše pogriješi, pa svaka od njih ima svoj test.
\end{frame}

% ------------------------------------------------------------
\begin{frame}{Matrica dopuštenih radnji}
Poslovna se pravila najkraće zapisuju tablicom stanja i radnji. Takva tablica ne postoji u kodu i ne da se iz njega izvesti: dolazi iz pravila domene, a napisana je prije koda, da se cjelina vidi na jednom mjestu.

\vspace{0.5em}
\begin{center}
{\small
\setlength{\tabcolsep}{4pt}
\renewcommand{\arraystretch}{1.4}
\begin{tabular}{@{}L{4.1cm}L{2.0cm}L{2.7cm}L{2.0cm}@{}}
  \toprule
  \textbf{Radnja} & \texttt{OPEN} & \texttt{PRELIM\_PASSED} & \texttt{CLOSED} \\
  \midrule
  nova prijava          & da & ne & ne \\
  promjena prijave      & da & da & ne \\
  odjava i ponovna prijava & da & da & ne \\
  čitanje popisa        & da & da & da \\
  \bottomrule
\end{tabular}
}
\end{center}

\vspace{0.6em}
Svaka ćelija s \enquote{ne} postaje jedna usporedba u sloju \codefile{services}. Korist je u tome što se rupe vide: ćelija na koju nitko nije pomislio ovdje upada u oči, a novo stanje dodaje stupac u kojem se za svaku radnju mora odlučiti.
\end{frame}

% ------------------------------------------------------------
\begin{frame}[fragile]{Pravilo pri novoj prijavi}
\begin{lstlisting}[language=python]
# app/services/registration_service.py
async def create_registration(db, comp_id, body, user):
    comp = await get_competition_or_404(db, comp_id)
    if get_competition_phase(comp) != Phase.OPEN:
        raise AppError("deadline_passed",
                       "Rok za nove prijave je istekao", 400)
    <@\kom{dohvat natjecatelja, vlasništvo, kategorija, upis}@>
    ...
\end{lstlisting}

\vspace{0.4em}
Provjera je prva naredba u funkciji, prije dohvata natjecatelja i prije provjere vlasništva. Zahtjev koji ne prolazi ne pokreće nijedan daljnji upit.

\vspace{0.4em}
Uvjet imenuje fazu u kojoj je radnja dopuštena, a ne faze u kojima se odbija. Ako se jednom doda četvrta faza, u njoj se radnja neće izvesti dok to netko izričito ne dopusti.
\end{frame}

% ------------------------------------------------------------
\begin{frame}[fragile]{Isto pravilo, drugi uvjet}
Izmjena postojeće prijave dopuštena je i nakon preliminarnog roka, pa njezin uvjet imenuje jedino fazu u kojoj se radnja odbija:

\begin{lstlisting}[language=python]
# app/services/registration_service.py
async def _get_registration_for_change(db, comp_id, reg_id, user):
    comp = await get_competition_or_404(db, comp_id)
    if get_competition_phase(comp) == Phase.CLOSED:
        raise AppError("deadline_passed",
                       <@"Izmjene nisu moguće nakon konačnog roka"@>, 400)
    reg = await _get_registration_or_404(db, comp_id, reg_id)
    _check_lifter_ownership(reg.lifter, user)
    return reg
\end{lstlisting}

\vspace{0.3em}
Oba pravila pozivaju istu funkciju za fazu i vraćaju isti kod pogreške, a razlikuju se u jednoj usporedbi. Matrica je time prepisana u kod, redak po redak.
\end{frame}

% ------------------------------------------------------------
\begin{frame}[fragile]{Faza u odgovoru}
Pravila provodi poslužitelj, ali sučelje treba znati koje radnje uopće ponuditi. Zato odgovor o natjecanju nosi i izračunatu fazu:

\begin{lstlisting}[language=python]
# app/schemas/competition.py
class CompetitionResponse(BaseModel):
    prelim_deadline: datetime
    final_deadline: datetime
    @computed_field
    @property
    def phase(self) -> Phase:
        return get_competition_phase(self)
\end{lstlisting}

\vspace{0.15em}
\texttt{computed\_field} uvrštava u odgovor vrijednost koja nastaje dok se odgovor sastavlja. U bazi odgovarajućeg stupca nema, pa vrijednost ne može zaostati za vremenom.

\vspace{0.15em}
Sučelje nije zaštita. Zahtjev koji stigne izvan faze odbija sloj \codefile{services}, bez obzira na to što je sučelje prikazalo.
\end{frame}

% ------------------------------------------------------------
\begin{frame}[fragile]{Vježba: faza koja se mijenja bez upisa}
Natjecanje \emph{Zimski kup} ima preliminarni rok 14.02.2027. Prijava Ivana Perića glasi kao u prvoj vježbi:

\begin{lstlisting}
POST /competitions/{comp_id}/registrations
{"lifter_id": <oznaka natjecatelja>, "category": "83", "total": 500}
\end{lstlisting}

\vspace{0.3em}
\begin{enumerate}
  \itemsep0.35em
  \item Kao \texttt{admin@plreg.hr} pošaljite \texttt{GET /competitions} i zabilježite \texttt{phase}.
  \item \texttt{PATCH /competitions/\{id\}} s tijelom \texttt{\{"prelim\_deadline": "2026-11-01T23:59:00+01:00"\}}.
  \item Ponovite \texttt{GET /competitions} i pogledajte \texttt{phase}.
  \item Prijavite se kao \texttt{behemot@plreg.hr} i pošaljite prijavu iz okvira na početku vježbe.
  \item Vratite \texttt{prelim\_deadline} na 14.02.2027.
\end{enumerate}

\vspace{0.4em}
Koji ste \texttt{code} dobili u četvrtom koraku i koji je upis promijenio fazu?
\end{frame}

% ============================================================
\sekcija{Prijelaz stanja umjesto brisanja}{Oznaka stanja u zapisu, prijelaz imenovan u adresi, ponovljeni zahtjev, ponovna prijava nakon odjave.}

% ------------------------------------------------------------
\begin{frame}{Cijena brisanja zapisa}
Klub koji odjavljuje natjecatelja mogao bi obrisati redak prijave. To ima posljedice i izvan zahtjeva:

\begin{itemize}
  \itemsep0.35em
  \item organizator ne vidi koliko je prijava bilo i koliko ih je povučeno
  \item klub koji se predomisli prijavljuje iznova, pa prijava nosi novo vrijeme upisa
\end{itemize}

\vspace{0.5em}
Zato se zapis zadržava, a mijenja mu se oznaka stanja:

\vspace{0.3em}
\begin{center}
{\small
\setlength{\tabcolsep}{0pt}
\renewcommand{\arraystretch}{1.35}
\begin{tabular}{@{}L{3.4cm}L{7.6cm}@{}}
  \toprule
  \textbf{\texttt{status}} & \textbf{Značenje} \\
  \midrule
  \texttt{active}    & prijava vrijedi i natjecatelj je na popisu \\
  \texttt{withdrawn} & prijava je povučena, ali zapis o njoj ostaje \\
  \bottomrule
\end{tabular}
}
\end{center}

\vspace{0.4em}
Popis koji organizator šalje sucima nastaje filtriranjem po \texttt{status=active}, pa povučena prijava ostaje u bazi, a s popisa izlazi.
\end{frame}

% ------------------------------------------------------------
\begin{frame}{Pohranjeno i izvedeno stanje}
I natjecanje i prijava imaju stanje, ali ne na isti način:

\vspace{0.5em}
\begin{center}
{\small
\setlength{\tabcolsep}{0pt}
\renewcommand{\arraystretch}{1.5}
\begin{tabular}{@{}L{3.2cm}L{4.0cm}L{4.0cm}@{}}
  \toprule
  & \textbf{\texttt{status} prijave} & \textbf{\texttt{phase} natjecanja} \\
  \midrule
  Gdje stoji        & stupac u tablici & izračun iz dvaju rokova \\
  Mijenja ga        & zahtjev klijenta & protek vremena \\
  Dopuštene vrijednosti čuva & ograničenje u bazi & nabrojiv tip \texttt{Phase} \\
  U odgovoru je     & polje iz zapisa & \texttt{computed\_field} \\
  \bottomrule
\end{tabular}
}
\end{center}

\vspace{0.6em}
Pohranjeno stanje treba radnju koja ga mijenja. Izvedeno stanje takvu radnju ne treba, ali se ne može postaviti rukom.
\end{frame}

% ------------------------------------------------------------
\begin{frame}[fragile]{Prijelaz stanja u adresi (I)}
Odjava mijenja jedno polje zapisa, pa se nameće najkraći put: poslati novu vrijednost tog polja.

\begin{lstlisting}
PATCH /competitions/7/registrations/41
{"status": "withdrawn"}
\end{lstlisting}

\vspace{0.4em}
Takav zahtjev kaže kamo vrijednost ide, ali ne i odakle dolazi, pa nosi tri problema:

\begin{itemize}
  \itemsep0.3em
  \item odjava i ponovna prijava postaju isti zahtjev, pa dijele jedno pravilo i jedan kod pogreške
  \item \texttt{withdrawn} u \texttt{withdrawn} uredan je upis, pa poslužitelj ne razlikuje prvu odjavu od ponovljene
  \item kategorija i stanje stižu u istom tijelu, pa jedan zahtjev obavlja dvije nepovezane promjene
\end{itemize}
\end{frame}

% ------------------------------------------------------------
\begin{frame}[fragile]{Prijelaz stanja u adresi (II)}
Stanje nije svojstvo koje klijent postavlja, nego posljedica radnje. Zato shema izmjene polja \texttt{status} nema:

\begin{lstlisting}[language=python]
# app/schemas/registration.py
class RegistrationUpdate(BaseModel):
    category: str | None = None
    total: int | None = Field(default=None, ge=0)
\end{lstlisting}

\vspace{0.4em}
Svaka radnja dobiva svoju adresu, a s njom i svoja pravila i svoj kod pogreške:

\begin{lstlisting}
POST /competitions/{comp_id}/registrations/{reg_id}/withdraw
POST /competitions/{comp_id}/registrations/{reg_id}/reactivate
\end{lstlisting}
\end{frame}

% ------------------------------------------------------------
\begin{frame}{Prijelaz stanja u adresi (III)}
Zadnji dio adrese imenuje prijelaz. Zahtjev \texttt{POST .../withdraw} ne stvara novi zapis i ne vraća novu adresu: nad postojećim zapisom izvodi jedan imenovani prijelaz.

\vspace{0.9em}
\begin{konvencija}[Izbor, ne pravilo]
Nijedan standard ovaj oblik ne propisuje, a u praksi se sreću oba: \texttt{PATCH} s poljem stanja, kakav smo odbacili, i posebna adresa s glagolom. Adrese inače gradimo od imena resursa, pa je glagol u adresi odstupanje koje treba opravdati.
\end{konvencija}

\vspace{0.9em}
Adresa time imenuje što se događa, a ne koje se polje mijenja.
\end{frame}

% ------------------------------------------------------------
\begin{frame}[fragile]{Odjava u kodu}
Funkcija \texttt{\_get\_registration\_for\_change} daje prijavu nad kojom se smije raditi, pa odjavi ostaje jedino prijelaz stanja:

\begin{lstlisting}[language=python]
# app/services/registration_service.py
async def withdraw_registration(db, comp_id, reg_id, user):
    reg = await _get_registration_for_change(
        db, comp_id, reg_id, user
    )
    if reg.status == "withdrawn":
        raise AppError("already_withdrawn",
                       <@"Prijava je već povučena"@>, 400)
    reg.status = "withdrawn"
    await db.flush()
    return reg
\end{lstlisting}

\vspace{0.3em}
Zapis ostaje na mjestu, a mijenja se jedan stupac. Upis ulazi u istu transakciju koju je otvorio zahtjev.
\end{frame}

% ------------------------------------------------------------
\begin{frame}{Dvostruki zahtjev}
Korisnik koji ne vidi odgovor, jer je veza spora ili se stranica zaglavila, pritisnut će gumb još jednom. Drugi zahtjev stiže nad zapisom koji je već povučen.

\vspace{0.6em}
Dva su moguća odgovora, i oba se susreću u praksi:

\begin{itemize}
  \item \textbf{Tiho prihvatiti.} Stanje je već ono koje se tražilo, pa odgovor je \texttt{200} i ništa se ne mijenja. Radnja je time idempotentna: ponovljeni zahtjev ostavlja isto stanje.
  \item \textbf{Odbiti.} Drugi zahtjev dobiva \texttt{400} i kod \texttt{already\_withdrawn}, jer traži prijelaz koji iz trenutnog stanja nije moguć.
\end{itemize}

\end{frame}

% ------------------------------------------------------------
\begin{frame}{Cijena obaju odgovora}
Svaki se izbor u nekom slučaju pokaže kao loš:

\vspace{0.4em}
\begin{itemize}
  \itemsep0.4em
  \item \textbf{Tiho prihvaćanje.} Klub je odjavio prijavu, predomislio se i vratio je. Zahtjev iz zaglavljene kartice stigne poslije toga i odjavljuje je po drugi put. Pogreške nema, a natjecatelja s popisa nema.
  \item \textbf{Odbijanje.} Odjava je prošla, ali se odgovor izgubio. Ponovljeni zahtjev vraća \texttt{400}, pa sučelje prikazuje pogrešku za radnju koja je uspjela.
\end{itemize}

\vspace{0.4em}
Izbor ovisi o tome je li važno samo konačno stanje ili i trenutak u kojem je prijelaz izveden. U prvom je slučaju tiho prihvaćanje blaže, u drugom ponovljeni zahtjev treba odbiti.

\vspace{0.4em}
\begin{konvencija}[Izbor, ne pravilo]
U ovom projektu drugi zahtjev se odbija. Sučelje zato \texttt{already\_withdrawn} mora prikazati kao obavijest da je odjava već provedena, a ne kao pogrešku.
\end{konvencija}
\end{frame}

% ------------------------------------------------------------
\begin{frame}[fragile]{Ponovna prijava nakon odjave}
Klub se predomislio i želi natjecatelja vratiti na popis. Ako pošalje istu prijavu kao prvi put, dobiva \texttt{409}:

\begin{lstlisting}
POST /competitions/7/registrations
{"lifter_id": 31, "category": "83"}

409 {"code": "duplicate", "message": "..."}
\end{lstlisting}

\vspace{0.25em}
Redak prijave nije obrisan, pa ga ograničenje jedinstvenosti i dalje vidi. Ponovna prijava zato ima vlastitu radnju, \texttt{reactivate\_registration}:

\begin{lstlisting}[language=python]
# app/services/registration_service.py
if reg.status == "active":
    raise AppError("already_active", <@"Prijava je već aktivna"@>, 400)
reg.status = "active"
\end{lstlisting}

\vspace{0.2em}
Ponovna prijava prolazi kroz istu provjeru roka kao i odjava, pa nakon konačnog roka nije moguća.
\end{frame}

% ------------------------------------------------------------
\begin{frame}[fragile]{Zajednička pravila triju radnji}
Izmjena, odjava i ponovna prijava sada su tri zasebne radnje, svaka sa svojom adresom. Sve tri traže ista tri pravila: rok, postojanje prijave i vlasništvo.

\begin{lstlisting}[language=python]
# app/services/registration_service.py
async def update_registration(db, comp_id, reg_id, body, user):
    reg = await _get_registration_for_change(...)   # pozivatelj 1
async def withdraw_registration(db, comp_id, reg_id, user):
    reg = await _get_registration_for_change(...)   # pozivatelj 2
async def reactivate_registration(db, comp_id, reg_id, user):
    reg = await _get_registration_for_change(...)   # pozivatelj 3
\end{lstlisting}

\vspace{0.3em}
Pravila stoje na jednom mjestu, a svaka ih radnja dobiva pozivom iste funkcije; razlikuju se jedino u onome što slijedi nakon tog poziva. Napisana u svakoj funkciji posebno, u četvrtoj bi se mogla zaboraviti.
\end{frame}

% ------------------------------------------------------------
\begin{frame}{Vježba: odjava i ponovna prijava}
Prijavite se kao klub \texttt{behemot@plreg.hr}. Natjecanje \emph{Zimski kup} mora biti u fazi \texttt{OPEN}; ako nije, vratite mu rokove na 14.02.2027. i 07.03.2027.

\vspace{0.8em}
\begin{enumerate}
  \itemsep0.4em
  \item Prijavite natjecatelja koji još nije prijavljen, u kategoriji koja odgovara njegovu spolu (za \texttt{M} npr. \texttt{"83"}, za \texttt{F} npr. \texttt{"63"}), i zabilježite oznaku prijave.
  \item Pošaljite \texttt{POST .../withdraw} nad tom prijavom, pa isti zahtjev još jednom.
  \item Dohvatite popis prijava, prvo bez parametara, pa s \texttt{?status=withdrawn}.
  \item Pokušajte istog natjecatelja prijaviti iznova, tijelom iz prvog koraka.
  \item Pošaljite \texttt{POST .../reactivate} i ponovite treći korak.
\end{enumerate}

\vspace{0.8em}
Koji ste \texttt{code} dobili u drugom i u četvrtom koraku i koja ih provjera podiže?
\end{frame}

% ============================================================
\sekcija{Pogreška kao odgovor}{Domenska iznimka, jedno mjesto prevođenja, izvori pogreške, izbor statusnog koda.}

% ------------------------------------------------------------
\begin{frame}[fragile]{Domenska pogreška}
Iznimku \texttt{AppError} upotrebljavamo od prvog predavanja, a nismo je otvorili.

\vspace{0.3em}
Sam odgovor sloj \codefile{services} ne sastavlja: HTTP je posao sloja \codefile{routers}. Umjesto toga podiže iznimku koja nosi sve potrebno:

\begin{lstlisting}[language=python]
# app/core/errors.py
class AppError(Exception):
    def __init__(self, code: str, message: str,
                 status_code: int = 400):
        self.code = code
        self.message = message
        self.status_code = status_code
# app/services/registration_service.py
raise AppError("deadline_passed",
               "Rok za nove prijave je istekao", 400)
\end{lstlisting}

\vspace{0.2em}
\texttt{code} je stalan i namijenjen programu koji odgovor obrađuje, \texttt{message} čovjeku i smije se mijenjati. Zadani \texttt{status\_code} je \texttt{400}.
\end{frame}

% ------------------------------------------------------------
\begin{frame}[fragile]{Jedno mjesto prevođenja}
Bez zajedničkog rukovatelja (engl. \emph{handler}) svaka bi ruta morala hvatati iznimku i sama graditi odgovor. Prijevod zato stoji na jednom mjestu:

\begin{lstlisting}[language=python]
# app/core/errors.py
async def app_error_handler(_request, exc: AppError) -> JSONResponse:
    return JSONResponse(
        status_code=exc.status_code,
        content={"code": exc.code, "message": exc.message},
    )

# app/main.py
app.add_exception_handler(AppError, app_error_handler)
\end{lstlisting}

\vspace{0.3em}
Iznimka podignuta bilo gdje u obradi izlazi kroz ovaj rukovatelj, kao tijelo istoga oblika. Ruta o tome ne zna ništa.

\vspace{0.3em}
Iznimka koja nije \texttt{AppError} ovdje se ne hvata. Ako baza ne odgovori, \texttt{OperationalError} izlazi kao \texttt{500}: to nije odbijen zahtjev, nego kvar.
\end{frame}

% ------------------------------------------------------------
\begin{frame}[fragile]{Tri izvora pogreške}
Zahtjev može pasti i prije nego što dođe do pravila, pa aplikacija registrira tri rukovatelja:

\vspace{0.5em}
\begin{center}
{\small
\setlength{\tabcolsep}{3pt}
\renewcommand{\arraystretch}{1.45}
\begin{tabular}{@{}L{4.4cm}L{5.1cm}L{1.6cm}@{}}
  \toprule
  \textbf{Iznimka} & \textbf{Podiže ju} & \textbf{Kod} \\
  \midrule
  \texttt{AppError} & pravilo u \codefile{services} & iz iznimke \\
  \texttt{RequestValidationError} & shema u \codefile{schemas} & \texttt{422} \\
  \texttt{StarletteHTTPException} & framework, za nepostojeću rutu & iz iznimke \\
  \bottomrule
\end{tabular}
}
\end{center}

\vspace{0.35em}
Tijelo je u svim trima slučajevima istoga oblika, s poljima \texttt{code} i \texttt{message}. Odgovor na neispravno tijelo nosi i popis polja koja nisu prošla provjeru:

\begin{lstlisting}[language=json]
{"code": "validation_error",
 "message": "Neispravni podatci u zahtjevu",
 "errors": [{"field": "category", "message": "Nepoznata kategorija"}]}
\end{lstlisting}
\end{frame}

% ------------------------------------------------------------
\begin{frame}{Statusni kod odbijenog pravila}
\texttt{401} kaže da pošiljatelj nije prepoznat, a \texttt{403} da jest, ali na radnju nema pravo. Pravilo koje radnju odbija iako je s pošiljateljem sve u redu bira među ovima:

\vspace{0.6em}
\begin{center}
{\small
\setlength{\tabcolsep}{0pt}
\renewcommand{\arraystretch}{1.45}
\begin{tabular}{@{}L{1.1cm}L{5.2cm}L{4.9cm}@{}}
  \toprule
  \textbf{Kod} & \textbf{Značenje} & \textbf{Primjer} \\
  \midrule
  \texttt{400} & u ovom stanju nije moguća & rok je prošao \\
  \texttt{409} & sukob s postojećim zapisom & natjecatelj je već prijavljen \\
  \texttt{422} & sadržaj se ne može obraditi & kategorija nije na popisu \\
  \bottomrule
\end{tabular}
}
\end{center}

\vspace{0.6em}
Opisi tih kodova u standardu dopuštaju više od jednoga izbora, pa se za prekršeno poslovno pravilo u praksi sreću sva tri.
\end{frame}

% ------------------------------------------------------------
\begin{frame}{Izbor koda u projektu}
Izbor se donosi jednom, za cijeli projekt:

\vspace{0.9em}
\begin{konvencija}
U ovom projektu \texttt{409} nosi jedino sukob s postojećim zapisom, \texttt{422} jedino pogrešku u obliku podataka, a sve ostalo je \texttt{400}.
\end{konvencija}

\vspace{0.9em}
Klijent po kodu grana obradu: na \texttt{422} označava polje u obrascu, na \texttt{400} prikazuje poruku i ostavlja obrazac kakav je, a na \texttt{409} osvježava prikaz, jer zapis već postoji.

\vspace{1.1em}
Projekt koji isto pravilo odbija jednom s \texttt{400}, a drugdje s \texttt{409}, tjera klijenta da oba koda obrađuje jednako, čime kod prestaje nositi informaciju.
\end{frame}

% ------------------------------------------------------------
\begin{frame}{Kodovi pogrešaka u projektu (I)}
Zahtjev može biti odbijen i prije nego što dođe do pravila domene. Te su pogreške iste u svakoj aplikaciji:

\vspace{0.6em}
\begin{center}
{\small
\setlength{\tabcolsep}{3pt}
\renewcommand{\arraystretch}{1.4}
\begin{tabular}{@{}L{3.8cm}L{2.3cm}L{4.9cm}@{}}
  \toprule
  \textbf{\texttt{code}} & \textbf{\texttt{status\_code}} & \textbf{Kada} \\
  \midrule
  \texttt{validation\_error}    & \texttt{422} & tijelo ne odgovara shemi \\
  \texttt{invalid\_credentials} & \texttt{401} & prijava ili token ne vrijede \\
  \texttt{token\_expired}       & \texttt{401} & pristupni token je istekao \\
  \texttt{forbidden}            & \texttt{403} & uloga ili vlasništvo ne dopušta \\
  \texttt{not\_found}           & \texttt{404} & zapis ne postoji ili nije vidljiv \\
  \bottomrule
\end{tabular}
}
\end{center}

\vspace{0.6em}
Prva tri koda uvedena su s prijavom, četvrti s pravima pristupa, a peti dolazi uz svaki dohvat zapisa po oznaci.
\end{frame}

% ------------------------------------------------------------
\begin{frame}{Kodovi pogrešaka u projektu (II)}
Ostali kodovi dolaze iz pravila ove domene:

\vspace{0.5em}
\begin{center}
{\small
\setlength{\tabcolsep}{3pt}
\renewcommand{\arraystretch}{1.2}
\begin{tabular}{@{}L{3.8cm}L{2.3cm}L{4.9cm}@{}}
  \toprule
  \textbf{\texttt{code}} & \textbf{\texttt{status\_code}} & \textbf{Kada} \\
  \midrule
  \texttt{duplicate}            & \texttt{409} & zapis bi prekršio jedinstvenost \\
  \texttt{has\_registrations}   & \texttt{409} & natjecatelj ima prijave \\
  \texttt{deadline\_passed}     & \texttt{400} & radnja nije moguća u ovoj fazi \\
  \texttt{already\_withdrawn}   & \texttt{400} & prijava je već povučena \\
  \texttt{already\_active}      & \texttt{400} & prijava je već aktivna \\
  \texttt{invalid\_category}    & \texttt{400} & kategorija ne odgovara spolu \\
  \texttt{invalid\_deadlines}   & \texttt{400} & poredak rokova bio bi poremećen \\
  \bottomrule
\end{tabular}
}
\end{center}

\vspace{0.5em}
Novi kod uvodi se kad klijent na njega mora reagirati drukčije. Inače je dovoljno promijeniti \texttt{message}.
\end{frame}

% ------------------------------------------------------------
\begin{frame}[fragile]{Pogreška koju vraća baza}
Dvostruku prijavu odbija baza, ograničenjem jedinstvenosti nad oznakama natjecatelja i natjecanja. Svoju odluku javlja iznimkom, i to pri izvođenju upisa:

\begin{lstlisting}[language=python]
# app/services/registration_service.py
try:
    return await registration_repo.create(db, reg)
except IntegrityError:
    raise AppError("duplicate",
                   <@"Natjecatelj je već prijavljen"@>, 409)
\end{lstlisting}

\vspace{0.3em}
Bez ovog prevođenja klijent dobiva \texttt{500}, jer je \texttt{IntegrityError} za aplikaciju nepredviđeno stanje. Taj je kod uz to netočan: zahtjev je u sukobu sa stanjem, a to je \texttt{409}.

\vspace{0.3em}
Provjera prije upisa ne bi bila dovoljna, jer se stanje može promijeniti između provjere i upisa.
\end{frame}

% ------------------------------------------------------------
\begin{frame}[fragile]{Vježba: isti oblik, tri rukovatelja}
Kao klub \texttt{behemot@plreg.hr}, uz \emph{Zimski kup} u fazi \texttt{OPEN}, šaljete prijavu Ivana Perića, s oznakom natjecatelja iz prve vježbe.

\begin{lstlisting}
POST /competitions/{comp_id}/registrations
{"lifter_id": <oznaka natjecatelja>, "category": "...", "total": 500}
\end{lstlisting}

\vspace{0.2em}
Za svaki korak zabilježite statusni kod i polje \texttt{code}.

\begin{enumerate}
  \itemsep0.15em
  \item \texttt{GET /competitions/99999}
  \item \texttt{GET /competitionsss}, dakle ruta koje nema
  \item Pošaljite prijavu s kategorijom \texttt{"999"} i pogledajte polje \texttt{errors}
  \item Pošaljite prijavu s kategorijom \texttt{"83"} dvaput zaredom
  \item U \codefile{app/main.py} zakomentirajte redak koji registrira \texttt{app\_error\_handler}, ponovite četvrti korak, pa vratite redak
\end{enumerate}

\vspace{0.2em}
Koja su dva koraka dala isti \texttt{code} iz različitih rukovatelja, a što je klijent dobio u petom?
\end{frame}

% ============================================================
\sekcija{Testiranje pravila koja ovise o vremenu}{Pravilo koje se ne može dočekati, zamrznuto vrijeme, zadano vrijeme datoteke, test bez baze.}

% ------------------------------------------------------------
\begin{frame}{Pravilo koje se ne da isprobati}
Pravilo o roku vrijedi samo u određenom razdoblju. Test koji ga želi provjeriti morao bi dočekati taj trenutak, a do roka su tjedni.

\vspace{0.6em}
Dva su načina da se čekanje izbjegne:

\begin{itemize}
  \item \textbf{Pomaknuti podatke.} Natjecanju se u testu postave rokovi tik uz sadašnji trenutak. Test prolazi, ali ishod ovisi o tome koliko se dugo izvodi, pa povremeno pada bez stvarnog razloga.
  \item \textbf{Pomaknuti vrijeme.} Trenutno vrijeme se za potrebe testa zamrzne na zadanu vrijednost, a podatci ostaju kakvi jesu.
\end{itemize}

\vspace{0.6em}
Drugi je put pouzdaniji, jer test ne ovisi ni o trajanju ni o danu kad se pokreće.
\end{frame}

% ------------------------------------------------------------
\begin{frame}[fragile]{Zamrznuto vrijeme}
\begin{lstlisting}[language=python]
# tests/test_registrations.py
@freeze_time("2026-04-20 10:00:00")
async def test_create_registration_prelim_passed(
    client, lifter, competition
):
    headers = await auth_header(client, "testclub", "klub123")
    resp = await client.post(
        url(competition),
        json={"lifter_id": lifter.id, "category": "83", "total": 500},
        headers=headers,
    )
    assert resp.status_code == 400
    assert resp.json()["code"] == "deadline_passed"
\end{lstlisting}

\vspace{0.2em}
Natjecanje iz \emph{fixturea} ima preliminarni rok 15.04.2026. i konačni 01.05.2026. Zamrznuto vrijeme 20.04. pada između njih, pa je faza \texttt{PRELIM\_PASSED} i nova prijava mora biti odbijena.
\end{frame}

% ------------------------------------------------------------
\begin{frame}[fragile]{Zadano vrijeme cijele datoteke}
Većina testova nad prijavama ne provjerava rok, ali svaki od njih o roku ovisi: bez zamrznutog vremena ishod bi im se mijenjao s danom pokretanja. Datoteka zato ima zadano vrijeme:

\begin{lstlisting}[language=python]
# tests/test_registrations.py
@pytest.fixture(autouse=True)
def open_phase():
    with freeze_time("2026-04-01 10:00:00"):
        yield
\end{lstlisting}

\vspace{0.4em}
\texttt{autouse=True} primjenjuje \emph{fixture} na svaki test u datoteci, i kad ga test ne navede u svojem potpisu. Vrijeme 01.04. pada prije preliminarnog roka, pa su sve radnje dopuštene.

\vspace{0.4em}
Test koji provjerava samo pravilo roka postavlja vlastito vrijeme oznakom \texttt{@freeze\_time}, koja ima prednost pred zadanim.
\end{frame}

% ------------------------------------------------------------
\begin{frame}[fragile]{Test funkcije bez baze}
Funkcija koja fazu računa ne dira bazu, pa je njezin test ne treba. Natjecanje se stvara u memoriji, a funkcija se poziva izravno:

\begin{lstlisting}[language=python]
# tests/test_phases.py
COMPETITION = Competition(
    name="Test Kup", location="Zadar", date=date(2026, 5, 16),
    prelim_deadline=datetime(2026, 4, 15, tzinfo=timezone.utc),
    final_deadline=datetime(2026, 5, 1, tzinfo=timezone.utc),
)

@freeze_time("2026-04-15 00:00:00")
def test_prelim_passed_exactly_at_prelim_deadline():
    assert get_competition_phase(COMPETITION) == Phase.PRELIM_PASSED
\end{lstlisting}

\vspace{0.2em}
Zapis modela stvoren u memoriji nije upisan u bazu, a funkciji to ni ne treba: ona čita dva polja. Test zato ne otvara transakciju i ne šalje zahtjev, nego poziva funkciju i gleda povratnu vrijednost.
\end{frame}

% ------------------------------------------------------------
\begin{frame}{Dvije tvrdnje u testu}
Test tvrdi dvoje: da je zahtjev odbijen i da je odbijen iz očekivanog razloga. Prvu tvrdnju nosi statusni kod, drugu kod pogreške.

\vspace{0.5em}
Bez druge tvrdnje test bi prošao i da je zahtjev pao zbog neispravne kategorije ili nepostojećeg natjecatelja, dakle iz razloga koji s rokom nema veze.

\vspace{0.6em}
Za svako pravilo potrebna su najmanje dva testa: jedan u kojem radnja prolazi i jedan u kojem je odbijena. Prvi pokazuje da pravilo ne zaustavlja ono što bi trebalo dopustiti.

\vspace{0.4em}
Pravilo koje razlikuje tri faze traži tri takva testa, po jedan za svaku fazu.
\end{frame}

% ------------------------------------------------------------
\begin{frame}[fragile]{Vježba: pomaknuta granica}
Testovi se pokreću iz mape \texttt{api/}, uz pokrenutu bazu:

\begin{lstlisting}[language=bash]
pytest tests/test_phases.py -v
\end{lstlisting}

\vspace{0.8em}
\begin{enumerate}
  \itemsep0.4em
  \item Pokrenite testove i pročitajte njihova imena: svako opisuje jednu granicu.
  \item U \codefile{app/core/phases.py} zamijenite \texttt{now < prelim} s \texttt{now <= prelim} i pokrenite ih ponovno.
  \item Pročitajte ispis testa koji je pao: koje je stanje očekivao, a koje dobio.
  \item Vratite uvjet na staro i provjerite da svi testovi prolaze.
\end{enumerate}

\vspace{0.8em}
Koji je test pao i koje značenje roka opisuje njegovo ime?
\end{frame}

% ------------------------------------------------------------
\begin{frame}{Sažetak}
\begin{itemize}
  \item Provjera ulaza gleda poslani sadržaj, poslovno pravilo stanje aplikacije.
  \item Pravilo stoji u sloju \codefile{services}, ispred svakog upisa u bazu.
  \item Shema pravilo provodi kad su svi podatci u tijelu zahtjeva; inače ga provodi sloj pravila.
  \item Izmjena, odjava i ponovna prijava dijele ista tri pravila, u istoj funkciji.
  \item Faza se računa iz dvaju rokova pri svakom zahtjevu, pa ne zaostaje za vremenom.
  \item Rok nosi oznaku vremenske zone u shemi, u stupcu i u usporedbi.
  \item Odjava mijenja oznaku stanja, a prijelaz stanja dobiva vlastitu adresu.
  \item Domensku pogrešku s kodom i porukom u odgovor prevodi jedan rukovatelj.
  \item Jedinstvenost provodi baza, a njezinu iznimku prevodimo u \texttt{409}.
  \item Pravilo koje ovisi o vremenu testira se zamrzavanjem vremena.
\end{itemize}
\end{frame}

\end{document}
